a specialization homomorphism:
\sgaStableTarget{sga6.e10.eq.7-14-1}@@M0@@
by defining, for every prime cycle @@M1@@ of @@M2@@ purely of dimension @@M3@@, identified with an integral subscheme of @@M4@@, the element @@M5@@ of @@M6@@ as the cycle associated with the closed subscheme @@M7@@ of @@M8@@, where @@M9@@ is the schematic closure of @@M10@@ in @@M11@@ (which is a closed subscheme of @@M12@@ flat over @@M13@@ (EGA IV 2.8), with generic fiber equal to the original @@M14@@, and satisfying @@M15@@ (EGA IV 7.1.13)). When @@M16@@ is locally equidimensional (i.e., all its irreducible components through a given point have the same dimension), one immediately verifies that the specialization homomorphism just defined induces homomorphisms
\sgaStableTarget{sga6.e10.eq.7-14-1bis}@@M17@@

Now suppose that @@M18@@ is quasi-projective, so that the notion of rational equivalence of cycles is available on @@M19@@ and on @@M20@@ (cf. [5], where the theory is developed over an algebraically closed field). It seems very plausible that homomorphism (7.14.1) carries cycles rationally equivalent to zero to cycles rationally equivalent to zero, and thus induces a corresponding homomorphism on the ``Chow'' groups of cycle classes
\sgaStableTarget{sga6.e10.eq.7-14-2}@@M21@@
or again, when @@M22@@ is locally equidimensional:
\sgaStableTarget{sga6.e10.eq.7-14-2bis}@@M23@@

One might try to prove this by a method paralleling that of 7.12, but in which the recourse to resolution of singularities and the reduction
